\section*{अध्याय \ref{ch:g11:functional-groups} --- क्रियात्मक समूह और परिवार}

\begin{solution}{exo:g11:functional-groups:1}
हाइड्रॉक्सिल, \omterm{def:g11:functional-groups:alcohol}{ऐल्कोहॉल}; श्रृंखला के भीतर कार्बोनिल, \omterm{def:g11:functional-groups:carbonyl}{कीटोन}; कार्बोक्सिल,
\omterm{def:g11:functional-groups:carboxylic-acid}{कार्बोक्सिलिक अम्ल}; \omterm{def:g11:functional-groups:carboxylic-acid}{एस्टर} समूह, \omterm{def:g11:functional-groups:carboxylic-acid}{एस्टर}; ऐमीनो समूह \ce{-NH2}, \omterm{def:g11:functional-groups:amine}{ऐमीन}।
\end{solution}

\begin{solution}{exo:g11:functional-groups:2}
मेथेनॉल; प्रोपेन-1-ऑल; मेथेनोइक अम्ल; प्रोपेनैल; प्रोपेनोन।
\end{solution}

\begin{solution}{exo:g11:functional-groups:3}
\omterm{def:g11:functional-groups:alcohol}{ऐल्कोहॉल}; \omterm{def:g11:functional-groups:carbonyl}{ऐल्डिहाइड}; \omterm{def:g11:functional-groups:carboxylic-acid}{एस्टर}; \omterm{def:g11:functional-groups:amine}{ऐमाइड}; \omterm{def:g11:functional-groups:halogenoalkane}{हैलोऐल्केन}; \omterm{def:g11:functional-groups:halogenoalkane}{ऐल्कीन}।
\end{solution}

\begin{solution}{exo:g11:functional-groups:4}
प्रोपेनैल \ce{CH3-CH2-CHO}, प्रोपेनोन \ce{CH3-CO-CH3}। प्रोपेनैल
\omterm{def:g11:functional-groups:carbonyl}{ऐल्डिहाइड} है: उसका \ce{C=O} वाला कार्बन श्रृंखला के सिरे पर है, एक हाइड्रोजन
\omterm{def:g7:atoms-and-molecules:atom}{परमाणु} से आबंधित; प्रोपेनोन में वह दो कार्बन \omterm{def:g7:atoms-and-molecules:atom}{परमाणुओं} के बीच है।
\end{solution}

\begin{solution}{exo:g11:functional-groups:5}
प्राथमिक (\ce{-OH} वाले कार्बन का एक कार्बन पड़ोसी है);
द्वितीयक (दो); तृतीयक (तीन)।
\end{solution}

\begin{solution}{exo:g11:functional-groups:6}
\setchemfig{atom sep=1.5em}
पेंटेन-3-ऑल का कंकाल \chemfig{-[:30]-[:-30](-[:-90]OH)-[:30]-[:-30]};
2-मेथिलब्यूटेनोइक अम्ल का कंकाल \chemfig{-[:30]-[:-30](-[:-90])-[:30](=[:90]O)-[:-30]OH};
हेक्सेन-2-ओन का कंकाल \chemfig{-[:30](=[:90]O)-[:-30]-[:30]-[:-30]-[:30]};
और प्रोपिल एथेनोएट का कंकाल \chemfig{-[:30](=[:90]O)-[:-30]O-[:30]-[:-30]-[:30]}।
\end{solution}

\begin{solution}{exo:g11:functional-groups:7}
\ce{C3H6O}। हाँ, वे \omterm{def:g11:organic-skeletons:isomer}{समावयवी} हैं, पर अलग समूहों के साथ: एक
\omterm{def:g11:functional-groups:carbonyl}{ऐल्डिहाइड} और एक \omterm{def:g11:functional-groups:carbonyl}{कीटोन} (दोनों कार्बोनिल समूह, अलग जगहों पर)।
\end{solution}

\begin{solution}{exo:g11:functional-groups:8}
\ce{HCOO-CH2-CH3}, एथिल मेथेनोएट; \ce{CH3-COO-CH2-CH2-CH3}, प्रोपिल
एथेनोएट।
\end{solution}

\begin{solution}{exo:g11:functional-groups:9}
मेथिल एथेनोएट (एथेनोइक अम्ल, मेथेनॉल); एथिल प्रोपेनोएट (प्रोपेनोइक
अम्ल, एथेनॉल); मेथिल मेथेनोएट (मेथेनोइक अम्ल, मेथेनॉल)।
\end{solution}

\begin{solution}{exo:g11:functional-groups:10}
\setchemfig{atom sep=1.5em}
\begin{center}
\small
\begin{tabular}{cccc}
\chemfig{-[:30]-[:-30]-[:30]-[:-30]OH} &
\chemfig{-[:30](-[:90]OH)-[:-30]-[:30]} &
\chemfig{-[:30](-[:90])-[:-30]-[:30]OH} &
\chemfig{-[:0](-[:90]OH)(-[:-90])-[:0]} \\
ब्यूटेन-1-ऑल & ब्यूटेन-2-ऑल & 2-मेथिलप्रोपेन-1-ऑल & 2-मेथिलप्रोपेन-2-ऑल \\
प्राथमिक & द्वितीयक & प्राथमिक & तृतीयक
\end{tabular}
\end{center}
\end{solution}

\begin{solution}{exo:g11:functional-groups:11}
एक \omterm{def:g11:functional-groups:amine}{ऐमाइड} (एथेनेमाइड)। \emph{-ईन}। \omterm{def:g11:functional-groups:carboxylic-acid}{कार्बोक्सिलिक अम्ल} (एक ही कार्बन पर \ce{C=O} और
\ce{-OH}), \omterm{def:g11:functional-groups:carboxylic-acid}{एस्टर} (\ce{C=O} और \ce{-O-}) और \omterm{def:g11:functional-groups:amine}{ऐमाइड}
(\ce{C=O} और \ce{N}): तीन परिवार।
\end{solution}

\begin{solution}{exo:g11:functional-groups:12}
ऐस्पिरिन: एक कार्बोक्सिल समूह (\omterm{def:g11:functional-groups:carboxylic-acid}{कार्बोक्सिलिक अम्ल}) और एक \omterm{def:g11:functional-groups:carboxylic-acid}{एस्टर} समूह।
पैरासिटामॉल: वलय पर एक फ़ीनॉल \ce{-OH} और एक \omterm{def:g11:functional-groups:amine}{ऐमाइड} समूह
\ce{NH-CO}।
\end{solution}

\begin{solution}{exo:g11:functional-groups:13}
\ce{C2H6O}। केवल एथेनॉल, अपने \ce{O-H} के साथ, अपने \omterm{def:g7:atoms-and-molecules:molecule}{अणुओं} के बीच \omterm{def:g11:polarity-and-cohesion:hydrogen-bond}{हाइड्रोजन
आबंध} बनाता है (मेथॉक्सीमेथेन के ऑक्सीजन पर कोई हाइड्रोजन नहीं है)।
एथेनॉल ऊँचे ताप पर, \qty{78}{\celsius} पर, उबलता है; मेथॉक्सीमेथेन
\qty{-25}{\celsius} पर।
\end{solution}

\begin{solution}{exo:g11:functional-groups:14}
एक हाइड्रॉक्सिल समूह और एक कार्बोक्सिल समूह। श्रृंखला में तीन कार्बन हैं,
अम्ल वाला कार्बन कार्बन 1 है और \ce{-OH} कार्बन 2 पर:
2-हाइड्रॉक्सीप्रोपेनोइक अम्ल। ग्लाइसीन: 2-ऐमीनोएथेनोइक अम्ल (अक्सर
ऐमीनोएथेनोइक अम्ल लिखा जाता है)।
\end{solution}

\begin{solution}{exo:g11:functional-groups:15}
एथेनॉल \ce{CH3-CH2-OH} है, \omterm{def:g11:organic-skeletons:formulas}{आण्विक सूत्र} \ce{C2H6O}, एक \omterm{def:g11:functional-groups:alcohol}{ऐल्कोहॉल}; एथेनैल \ce{CH3-CHO} है,
\omterm{def:g11:organic-skeletons:formulas}{आण्विक सूत्र} \ce{C2H4O}, एक \omterm{def:g11:functional-groups:carbonyl}{ऐल्डिहाइड}; एथेनोइक अम्ल \ce{CH3-COOH} है, \omterm{def:g11:organic-skeletons:formulas}{आण्विक सूत्र} \ce{C2H4O2},
एक \omterm{def:g11:functional-groups:carboxylic-acid}{कार्बोक्सिलिक अम्ल}। एथेनॉल से एथेनैल तक दो हाइड्रोजन \omterm{def:g7:atoms-and-molecules:atom}{परमाणु} खोते हैं;
एथेनैल से एथेनोइक अम्ल तक एक ऑक्सीजन \omterm{def:g7:atoms-and-molecules:atom}{परमाणु} जुड़ता है।
\end{solution}

\begin{solution}{pb:g11:functional-groups:1}
\textbf{1.} A: \omterm{def:g11:functional-groups:carboxylic-acid}{एस्टर} समूह; B: हाइड्रॉक्सिल; C: कार्बोक्सिल; D: श्रृंखला के
सिरे पर कार्बोनिल; E: \omterm{def:g11:functional-groups:carboxylic-acid}{एस्टर} समूह।

\textbf{2.} A \omterm{def:g11:functional-groups:carboxylic-acid}{एस्टर}, B \omterm{def:g11:functional-groups:alcohol}{ऐल्कोहॉल}, C \omterm{def:g11:functional-groups:carboxylic-acid}{कार्बोक्सिलिक अम्ल}, D \omterm{def:g11:functional-groups:carbonyl}{ऐल्डिहाइड}, E \omterm{def:g11:functional-groups:carboxylic-acid}{एस्टर}।

\textbf{3.} A (केला) और E (अनानास): फलों जैसी गंधें।

\textbf{4.} प्राथमिक: \ce{-OH} वाला कार्बन एक कार्बन से आबंधित है।

\textbf{5.} \ce{C6H12O}; उदाहरण के लिए हेक्सेन-2-ओन,
\ce{CH3-CO-CH2-CH2-CH2-CH3}।

\textbf{6.} \ce{C-OH} वाली श्रृंखला में चार कार्बन हैं, जिन पर
\ce{-OH} की ओर से क्रमांक दिए गए हैं, और कार्बन 3 पर एक मेथिल है: 3-मेथिलब्यूटेन-1-ऑल।

\textbf{7.} C: एथेनोइक अम्ल; D: हेक्सेनैल।

\textbf{8.} एथिल ब्यूटेनोएट, ब्यूटेनोइक अम्ल और एथेनॉल से संबंधित।

\textbf{9.} \emph{एथेनोएट}: अम्ल भाग \ce{CH3-CO}, एथेनोइक
अम्ल से; \emph{3-मेथिलब्यूटिल}: \omterm{def:g11:functional-groups:alcohol}{ऐल्कोहॉल} भाग, 3-मेथिलब्यूटेन-1-ऑल से।

\textbf{10.} C (एथेनोइक अम्ल) और B (3-मेथिलब्यूटेन-1-ऑल)।

\textbf{11.} A \ce{C7H14O2}; B \ce{C5H12O}; C \ce{C2H4O2}।

\textbf{12.} $\ce{C2H4O2} + \ce{C5H12O}$ में \ce{C7H16O3} है, यानी एक
\ce{H2O} और \ce{C7H14O2}: खोया हुआ \omterm{def:g7:atoms-and-molecules:molecule}{अणु} पानी है।

\textbf{13.} \ce{C2H4O2 + C5H12O -> C7H14O2 + H2O}: C 7 = 7, H 16 = 16,
O 3 = 3.

\textbf{14.} $M(\text{B}) = 5 \times 12.0 + 12 \times 1.0 + 16.0 =
\qty{88.0}{g/mol}$; $M(\text{C}) = \qty{60.0}{g/mol}$।

\textbf{15.} $60.0 + 88.0 - 18.0 = \qty{130.0}{g/mol}$।

\textbf{16.} $7 \times 12.0 + 14 \times 1.0 + 2 \times 16.0 =
\qty{130.0}{g/mol}$: वही।

\textbf{17.} E, \ce{C6H12O2}: $6 \times 12.0 + 12 \times 1.0 + 2 \times
16.0 = \qty{116.0}{g/mol}$।

\textbf{18.} \ce{C7H14O2}: A वाला ही सूत्र, इसलिए \omterm{def:g11:organic-skeletons:isomer}{समावयवी}, और
वही \omterm{def:g10:the-mole:molar-mass}{मोलर द्रव्यमान}। केवल \omterm{def:g10:the-mole:molar-mass}{मोलर द्रव्यमान} \omterm{def:g11:organic-skeletons:isomer}{समावयवियों} में भेद नहीं कर सकता।

\textbf{19.} A और हेप्टेनोइक अम्ल में वही \omterm{def:g7:atoms-and-molecules:atom}{परमाणु} हैं पर अलग
समूह, एक \omterm{def:g11:functional-groups:carboxylic-acid}{एस्टर} और एक \omterm{def:g11:functional-groups:carboxylic-acid}{कार्बोक्सिलिक अम्ल}: एक की गंध केले जैसी है, दूसरे की
बासी। समूह, और वह कहाँ लगा है, यही तय करता है।

\textbf{20.} \qty{130.0}{g/mol}।
\end{solution}
